% ============================================================================
% CAPÍTULO 6 — CONCLUSÃO
% ============================================================================
\chapter{Conclusão}
\label{cap:conclusao}

A pergunta de pesquisa que orientou este trabalho (Capítulo
\ref{cap:introducao}) era se níveis induzidos de personalidade Big Five em
agentes de inteligência artificial predizem o escore do questionário de
perfil de investidor (IPA). A resposta, com os dados da análise principal
(Experimento A, BFI-44, N=300, \texttt{llama3.1:8b}), é afirmativa, porém
parcial: o modelo com os cinco níveis induzidos explica 45,9\% da
variância do escore IPA ($R^2=0{,}459$; $p<0{,}001$; Tabela \ref{tab:ols}),
mas apenas dois dos cinco traços, Abertura à Experiência e Extroversão,
têm efeito que se sustenta em todas as rodadas da análise de robustez e
generalização. As três questões de pesquisa definidas na introdução
detalham essa resposta.

\section{Respostas às questões de pesquisa}

\textbf{RQ1: Os LLMs conseguem simular consistentemente os traços de
personalidade solicitados?} Sim, mas não em qualquer modelo. No
Experimento A, a correlação entre o nível induzido e o autorrelato no
BFI-44 ficou entre 0,813 e 0,929 nos cinco traços (Tabela
\ref{tab:manipulacao}), com valores semelhantes quando cada traço foi
induzido isoladamente (Experimento B, Tabela \ref{tab:manipulacao-exp-b}),
em todas as temperaturas testadas e nos demais modelos da análise de
generalização, de 7 a 12 bilhões de parâmetros (seção
\ref{sec:resultados-robustez}). A indução também
se mostrou específica: cada nível induzido se reflete sobretudo no
próprio traço, e a correlação entre os traços autorrelatados foi baixa
(média de 0,110 em módulo; seção \ref{sec:resultados-intercorrelacao}).
Essa consistência, porém, não é garantida em qualquer modelo: o modelo
mais leve testado na seleção do modelo (\texttt{llama3.2:1b}) não
incorporou nenhum dos cinco traços de forma detectável, e o
\texttt{llama3.2:3b} incorporou apenas parte deles (seção
\ref{sec:selecao-modelo}). Dentro da família Llama, a fidelidade cresceu
com o tamanho do modelo, mas o \texttt{gemma3:4b}, com 4,3 bilhões de
parâmetros, também incorporou os cinco traços no teste reduzido (N=30),
o que indica que a capacidade de simular a personalidade depende do
modelo, e não apenas do seu porte.

\textbf{RQ2: Quais associações entre personalidade induzida e perfil de
investidor são robustas?} As de \textbf{Abertura à Experiência} e
\textbf{Extroversão}, ambas positivas: níveis induzidos mais altos desses
traços associam-se a escores IPA mais altos (mais arrojados). As duas
foram significativas, com $p<0{,}001$, nas oito rodadas com 300 agentes
da análise de robustez e generalização (o Experimento A e mais sete
rodadas: três temperaturas adicionais, um novo sorteio de perfis e três
outros modelos abertos; Tabela \ref{tab:robustez-sintese}) e aparecem na
mesma direção nos dois modelos Nemotron. Abertura à Experiência é, além
disso, o único traço com efeito significativo que converge com os dois
artigos-base
\cite{debortoli2019,jiangpengyan2024} e com o estudo de
\citeonline{hartley2025}, e o único que se manteve significativo quando
induzido isoladamente (Experimento B). Os efeitos de
\textbf{Consciência} e \textbf{Neuroticismo}, significativos no
Experimento A, são instáveis: perdem significância com outro
sorteio de perfis e mudam de significância ou de sinal entre modelos;
com a temperatura, variam em intensidade. \textbf{Amabilidade} não tem
efeito consistente.

\textbf{RQ3: Em que medida essas associações e os perfis financeiros
dependem do modelo e dos parâmetros de geração?} A \textit{direção} das
associações robustas depende pouco: Abertura à Experiência e Extroversão
mantêm o sinal e a significância em todas as temperaturas e em todos os
modelos da análise de generalização. A temperatura afeta principalmente o
ruído: reduzi-la de 0,7 para 0,0 eleva o $R^2$ de 0,459 para 0,707 e torna
a manipulação mais fiel, sem mudar o sinal e a significância de Abertura
à Experiência e de Extroversão; Amabilidade, porém, torna-se
significativa nas temperaturas mais baixas, o que mostra que uma única
configuração de geração pode indicar efeitos que não se sustentam. Já o
\textit{tamanho} dos efeitos e, sobretudo, o \textit{perfil de investidor
absoluto} dependem fortemente do modelo: com os mesmos perfis de
personalidade, a proporção de agentes Arrojados vai de 14,0\%
(\texttt{llama3.1:8b}) a 91,5\% (\texttt{mistral:7b}), e o perfil
atribuído a um mesmo agente coincide com o obtido no \texttt{llama3.1:8b}
em 20,3\% a 62,0\% dos casos, conforme o modelo. Os efeitos instáveis
(Consciência e Neuroticismo), assim como o de Amabilidade, são justamente
os que mais dependem do modelo e da geração.

Esse padrão de resultado (dois efeitos estáveis cercados de efeitos
menores que não se sustentam) delimita o que este tipo de simulação pode
oferecer. A convergência de Abertura à Experiência entre um estudo com
124 estudantes, um estudo com 3.325 investidores reais e a simulação com
agentes de IA deste trabalho é a evidência mais forte que este trabalho
produz em favor de uma resposta afirmativa à pergunta de pesquisa. Essa
convergência, contudo, não demonstra que um agente de linguagem
``possua'' personalidade no sentido psicológico do termo; o que ela
mostra é que, com indução e checagem de manipulação, a simulação
reproduziu a direção de uma relação relatada nos dois estudos com
humanos.

\section{Objetivos específicos}

Retomando os objetivos específicos definidos na introdução, um a um:

\begin{enumerate}
    \item \textbf{Induzir e verificar níveis de personalidade por
    instrução}: atingido. A checagem de manipulação confirmou incorporação
    forte e estatisticamente significativa dos cinco traços tanto no
    Experimento A ($\rho$ entre 0,813 e 0,929, Tabela \ref{tab:manipulacao})
    quanto no Experimento B ($\rho$ entre 0,837 e 0,930 por traço, Tabela
    \ref{tab:manipulacao-exp-b}), com o \texttt{llama3.1:8b} escolhido por
    ter sido, na seleção do modelo (feita com o BFI-10), o único que
    reuniu manipulação significativa nos cinco traços (forte em quatro
    deles e moderada em Extroversão), N suficiente para inferência e
    variância adequada no escore do IPA;
    \item \textbf{Aplicar o IPA e obter escore contínuo e classificação
    categórica}: atingido. Os 300 agentes do Experimento A foram
    classificados, pelo escore total do IPA (8 perguntas, 8 a 32 pontos,
    seção \ref{sec:ipa-pontuacao}), em 3,7\% Conservador, 82,3\%
    Moderado e 14,0\% Arrojado (Tabela \ref{tab:distribuicao-ipa}); para a
    análise estatística, foi usado o escore contínuo restrito às
    perguntas 2 a 8 (média 17,58, desvio padrão 2,72), que exclui a
    pergunta 1, cuja resposta foi sorteada;
    \item \textbf{Testar a associação estatística entre nível induzido e
    escore do IPA}: atingido para os cinco traços, por correlação de
    Spearman (Tabela \ref{tab:spearman-ipa}) e regressão OLS (Tabela
    \ref{tab:ols}), com quatro dos cinco coeficientes estatisticamente
    significativos no Experimento A, dois deles estáveis (RQ2);
    \item \textbf{Comparar os achados com \citeonline{debortoli2019} e
    \citeonline{jiangpengyan2024}}: atingido (seção
    \ref{sec:discussao-alinhamento}). Abertura à Experiência converge com
    os dois artigos-base; Amabilidade converge com
    \citeonline{debortoli2019} e com a ausência de efeito sobre a alocação
    em \citeonline{jiangpengyan2024}; Extroversão diverge de
    \citeonline{debortoli2019}, mas é compatível com a menor aversão ao
    risco declarada relatada por \citeonline{jiangpengyan2024}; e
    Consciência e Neuroticismo, além de divergirem dos dois artigos-base,
    mostraram-se instáveis, o que impede tratar essa divergência como um
    achado substantivo;
    \item \textbf{Comparar diferentes modelos de linguagem quanto à
    qualidade da indução}: atingido (Capítulo \ref{cap:metodologia}, seção
    \ref{sec:selecao-modelo}, e Capítulo \ref{cap:resultados}, seção
    \ref{sec:resultados-robustez}). Na seleção do modelo, a qualidade da
    indução cresceu com o tamanho do modelo dentro da família Llama; na
    análise de generalização, todos os modelos de 7 bilhões de parâmetros
    ou mais incorporaram os cinco traços;
    \item \textbf{Testar se a indução simultânea dos cinco traços afeta a
    qualidade da indução ou a força da associação com o IPA}: atingido
    (Experimento B, seção \ref{sec:discussao-exp-b}). A indução simultânea
    não prejudica a checagem de manipulação em si, mas o desenho isolado,
    com N entre 49 e 69 por traço, não teve poder estatístico suficiente
    para confirmar ou refutar com segurança os efeitos secundários do
    Experimento A;
    \item \textbf{Avaliar a robustez e a generalização das associações}:
    atingido (seções \ref{sec:resultados-robustez} e
    \ref{sec:discussao-robustez}). Abertura à Experiência e Extroversão se
    mantiveram significativas em todas as temperaturas, sorteios e modelos
    da análise de robustez e generalização; Consciência e Neuroticismo,
    não.
\end{enumerate}

\section{Contribuições}

Este trabalho contribui em três frentes. Em primeiro lugar, produz
evidência de que agentes de IA baseados em LLM, quando adequadamente
induzidos e submetidos a uma checagem de manipulação, reproduzem
parcialmente a direção de relações já documentadas entre personalidade
Big Five e comportamento de risco financeiro: a convergência de Abertura
à Experiência entre os dois artigos-base e este trabalho é o achado
central nesse sentido, ao passo que a instabilidade dos demais traços
delimita com igual clareza até onde essa replicação se sustenta.

Em segundo lugar, o trabalho aplica e documenta uma metodologia replicável
para o uso de LLMs como sujeitos simulados em pesquisa comportamental:
indução de personalidade por instrução, checagem de manipulação
obrigatória antes de interpretar qualquer efeito sobre a variável
dependente e uma análise de robustez e generalização que repete o
experimento variando, um de cada vez, a temperatura, o sorteio de perfis
e o modelo. Essa análise mostrou que um efeito estatisticamente
significativo em uma única configuração (como Consciência e Neuroticismo
no Experimento A) pode não se sustentar em outras, e que, sem ela, este
trabalho teria apresentado quatro efeitos onde apenas dois são estáveis.
A comparação entre modelos na seleção do modelo também mostrou que
modelos menores (\texttt{llama3.2:1b} e \texttt{llama3.2:3b}) não
atingiram manipulação confiável dos cinco traços (Tabela
\ref{tab:comparacao-modelos}), o que recomenda verificar a manipulação
antes de usar um modelo nesse tipo de simulação.

Em terceiro lugar, a comparação entre modelos mostra que o modelo de
linguagem não é um detalhe de implementação: ele determina se a indução
funciona (o modelo mais leve testado não incorporou nenhum traço) e,
mesmo entre modelos em que a indução funciona, determina o perfil de
investidor absoluto atribuído a uma mesma personalidade. O perfil
absoluto obtido nesse tipo de simulação deve, portanto, ser reportado
como propriedade do modelo usado, e não como estimativa do comportamento
de investidores.

\section{Recomendações para trabalhos futuros}

\begin{itemize}
    \item Substituir o questionário de perfil de investidor por uma
    \textbf{tarefa comportamental de investimento}, ou complementá-lo com
    ela, em que o agente toma decisões (por exemplo, alocar um valor entre
    ativos com risco e retorno conhecidos) em vez de declarar
    preferências. Em \citeonline{borman2024}, personas de LLM
    representaram melhor o comportamento humano nesse formato do que em
    questionário, e essa mudança pode reduzir a instabilidade observada
    nos traços de efeito menor;
    \item Estender a análise de generalização a modelos maiores, inclusive
    proprietários, com N suficiente para inferência, já que os modelos
    Nemotron foram testados apenas em caráter exploratório (N=30);
    \item Adaptar o IPA, ou adotar outro instrumento de tolerância ao
    risco validado, ao contexto específico de agentes de IA, removendo ou
    ajustando itens originalmente pensados para respondentes humanos (como
    a pergunta de faixa etária, cuja resposta foi atribuída aleatoriamente
    neste trabalho por não fazer sentido que o modelo a respondesse; seção
    \ref{sec:limitacoes-pipeline}).
\end{itemize}
